\begin{lemma}[A matrix and a vector in its kernel]
\label{lem:selmer-local-incidence}
Let $(P,\mathfrak m,k)$ be a Noetherian regular local ring of dimension
$d$.  Let $Z=(z_{ij})$ be a $g\times g$ matrix and $y=(y_j)$ a column
vector, where $g\geq1$.  Suppose their entries lie in $\mathfrak m$ and
have linearly independent classes in $\mathfrak m/\mathfrak m^2$.
Put $I=((Zy)_1,\ldots,(Zy)_g)$ and $J=I+(\det Z)$.
Then $\sqrt J$ is prime and $\dim(P/\sqrt J)\leq d-g$.
If a prime $\mathfrak p\supseteq I$ avoids some $y_j$ and satisfies
$\dim(P/\mathfrak p)\geq d-g$, then
\[
                   \mathfrak p=\sqrt J.
\]
In particular, $\mathfrak p$ lies below every prime containing $J$.
\end{lemma}

\begin{proof}
Consider the closed subscheme $Y\subset\mathbf P^{g-1}_P$ parametrizing
lines $L$ such that $ZL=0$ and $y\in L$.  Its projection is proper and
has image $V(J)$: a nonzero $y$ determines the required line, while
$y=0$ permits any line in the kernel of the singular matrix $Z$.
On the chart where the $c$th coordinate of the line is $1$, its equations
are
\[
 z_{ic}+\sum_{j\ne c}u_jz_{ij}=0\quad(1\leq i\leq g),
 \qquad y_j-u_jy_c=0\quad(j\ne c).
\]
Every closed point lies above the closed point of $\operatorname{Spec}P$.
At such a point the ambient chart local ring is regular of dimension
$d+g-1$.  Its base parameters retain independent cotangent classes,
including when the point has a nontrivial residue field extension.
The displayed equations have independent linear parts: use the
coordinates $z_{ic}$ and $y_j$ with $j\ne c$ as pivots.  Thus the local
ring of $Y$ is regular of dimension $d-g$.

Regularity at all closed points implies regularity everywhere.  The
irreducible components of $Y$ are consequently disjoint.  Properness
forces every component to meet the closed fibre
$\mathbf P^{g-1}_k$, which is connected; hence $Y$ is irreducible.
Its proper image $V(J)$ is irreducible too, so $K=\sqrt J$ is prime.
Moreover, $\dim Y=d-g$ and a closed surjective map of Noetherian spaces
cannot increase dimension: lift a chain of irreducible closed sets by
successively choosing components of their inverse images.  Therefore
$\dim(P/K)\leq d-g$.

For the prime $\mathfrak p$ in the statement, the image of $y$ is a
nonzero kernel vector over $\operatorname{Frac}(P/\mathfrak p)$.
Thus $\det Z\in\mathfrak p$ and $K\subseteq\mathfrak p$.
A strict inclusion would give
$\dim(P/\mathfrak p)+1\leq\dim(P/K)\leq d-g$, a contradiction.
The regular-local facts and the test at closed points are recorded in
the Stacks Project, tags
\href{https://stacks.math.columbia.edu/tag/00NQ}{00NQ} and
\href{https://stacks.math.columbia.edu/tag/02IT}{02IT}.
\end{proof}

\paragraph{Use in the ordinary deformation argument.}
The reducible local presentation has a finite algebra $A=B\oplus N$,
with products of elements of $N$ lying in the maximal ideal of $B$.
Let $m_1,\ldots,m_g$ generate $N$ and let
$\mathfrak r\subset B$ be generated by these products.  The required
arithmetic input gives a surjection from a regular power series ring
$P$ of dimension $1+D+g$ to $A$, together with relations $Zm=0$ whose
matrix entries and vector coordinates have independent cotangent
classes.  The lemma then supplies the needed local prime inclusion.
Indeed, a prime $\mathfrak p\subset B$ with
$\dim(B/\mathfrak p)\geq1+D$ and
$\mathfrak r\nsubseteq\mathfrak p$ lifts to a prime avoiding some
$m_j$.  The determinant of the relation matrix annihilates $N$ and
therefore annihilates $\mathfrak r$.  It belongs to every prime
$\mathfrak q$ for which either $\mathfrak r\nsubseteq\mathfrak q$ or
$\operatorname{Ann}_B N\subseteq\mathfrak q$.  Apply the lemma to primes
above $\mathfrak p,\mathfrak q$ and contract to obtain
$\mathfrak p\subseteq\mathfrak q$.
The remaining arithmetic work is to establish this presentation and
its dimension bound for the actual deformation problem.
